\documentclass{article}
\usepackage{graphicx}
\usepackage{amsmath}
\usepackage{amssymb}
\usepackage{amsfonts}
\usepackage{hyperref}
\usepackage{exercises}
\usepackage{indentfirst}
\usepackage{circledsteps}
\usepackage{amsthm}
\usepackage{csquotes}
\usepackage{yhmath}
\usepackage{comment}
\usepackage{placeins}
\usepackage{geometry}
\usepackage{cleveref}
\usepackage{stmaryrd}
\usepackage{hyperref}
\usepackage{tikz}
\usepackage{tabto}
\newcommand*{\oldneg}{\mathord{\sim}}
\newcommand{\Z}{\mathbb{Z}}
\newcommand{\N}{\mathbb{N}}
\let\oldemptyset\emptyset
\newcommand{\lbc}{\llbracket}
\newcommand{\rbc}{\rrbracket}
\newcommand{\ltop}[1]{\stackrel{\leftarrow}{#1}}

\theoremstyle{definition}
\newtheorem{definition}{Definition}
\newtheorem{theorem}{Theorem}[section]
\newtheorem{corollary}{Corollary}[theorem]
\newtheorem{lemma}[theorem]{Lemma}

\setlength{\parindent}{.5in}
\geometry{margin=0.5in}
\setlength{\parskip}{12pt}
\title{Chapter 19 Notes}
\begin{document}
\date{}
\maketitle

\section*{Section 2}

\subsection*{Notes}
For the model in this section to work, we need $\omega_2\in M$. However, it is not clear to my why this is.

If we assume that $\omega_2\in M$, here is why $\langle\mathcal{G},\mathcal{R}\rangle$ is an $M$-set. It 
is easy to see that each particular forcing condition is in $M$ - each particular one can be explicitly defined, 
as it is assumed that $p(\alpha)=\langle\emptyset,\emptyset\rangle$ for all but a finite number of $\alpha\in\omega_2$. 
The set $\mathcal{G}$ of all forcing conditions is a subset of 
$\mathcal{P}(\mathcal{P}(\omega_2\times(\mathcal{P}(\omega\times\omega)\cap M))\cap M)\cap M$, which is in $M$.
$\mathcal{G}$ is also easily seen to be definable over $M$, and consequently, $\mathcal{G}\in M$. Then 
$\mathcal{R}$ is also easily seen to be in $M$, as it is definable over $M$, and is a subset of $\mathcal{G}\times\mathcal{G}$, which 
is in $M$ (as we have shown $\mathcal{G}$ to be in $M$).

\subsection*{Exercise 2.1}
This is basically the same proof as that of Theorem 4.1 of Chapter 18.

\begin{proof}

We prove that, for each $\alpha\in\omega_2$, $\lbc\oldneg\textrm{Const}(f_\alpha)\rbc$ is valid in the model
$M$. Suppose, to the contrary, that $p\nvDash\lbc\oldneg\textrm{Const}(f_\alpha)\rbc$ for some $p$. 
Then for some $p'$ with $p\mathcal{R}p'$, $p'\vDash\lbc\textrm{Const}(f_\alpha)\rbc$. By Theorem 3.2 of 
Chapter 18, it follows that $p\vDash\Diamond\lbc f_\alpha\approx\hat{c}\rbc$ for some contrsuctible set $c\in M$.
This means that, for some $r$ with $p'\mathcal{R}r$, $r\vDash\lbc f_\alpha\approx\hat{c}\rbc$. We derive a 
contradiction from this. $r$ is some forcing condition, so 
suppose that $r(\alpha)=\langle P,N\rangle$. There are then two cases.

Case 1) $c\subseteq P$. Because both $P$ and $N$ are finite, we can pick some integer $n$ 
such that $n\notin P\cup N$ (thus also $n\notin c$). Let $P'=P\cup\{n\}$. Then there is some 
forcing condition $q$ such that, for all $\beta\in\omega_2$, if $\beta\ne\alpha$, then 
$q(\beta)=r(\beta)$, and if $\beta=\alpha$, then $q(\beta)=\langle P',N\rangle$. Obviously 
$r\mathcal{R}q$. Because $c\in P'$, it follows that 
$q\vDash\lbc\hat{n}\in f_\alpha\rbc$. But because $n\notin c$, $q\vDash\lbc\oldneg(\hat{n}\in\hat{c})\rbc$. 
These two things jointly imply $q\nvDash\lbc f_\alpha\approx\hat{c}\rbc$. But 
also $q\vDash\lbc f_\alpha\approx\hat{c}\rbc$, as $r\vDash\lbc f_\alpha\approx\hat{c}\rbc$ and 
$r\mathcal{R}q$. Thus we have our contradiction in this case.

Case 2) $c\nsubseteq P$. Pick some $n$ in $c$ but not in $P$. Let $N'=N\cup\{n\}$. 
Then there is some forcing condition $q$ such that, for all $\beta\in\omega_2$, if 
$\beta\ne\alpha$, then $q(\beta)=r(\beta)$, and if $\beta=\alpha$, then $q(\beta)=\langle P,N'\rangle$. Obviously 
$r\mathcal{R}q$. Because $n\in c$, it follows that 
$q\vDash\lbc\hat{n}\in\hat{c}\rbc$. But because $n\in N'$, it follows that $q\vDash\lbc\oldneg(\hat{n}\in f_\alpha)\rbc$ (Reason: 
Suppose, for some $q'$ with $q\mathcal{R}q'$, that $q'\vDash\hat{n}\in f_\alpha$. Then for some 
$x$, $q'\vDash\lbc x\epsilon f_\alpha\rbc$ and $q'\vDash\lbc x\approx\hat{n}\rbc$. This means 
$x$ is $\hat{a}$ for some $a$, and thus, by Lemma 3.2 of Chapter 17, $a$ and $n$ are the same set. But 
this means $q'\vDash\lbc\hat{n}\epsilon f_\alpha\rbc$, which cannot happen, as $n\in N'$). These 
two things imply that $q\nvDash\lbc f_\alpha\approx\hat{c}\rbc$. But also $q\vDash\lbc f_\alpha\approx\hat{c}\rbc$ (for the 
same reason as in the first case), which is our contradiction.

\end{proof}

\end{document}
